\section{Introduction}
\label{sec:intro}

When the UK government changes a tax or benefit, two questions follow immediately. The first is static: who pays what the morning after? The second is dynamic: how do consumption, investment, incomes, and prices adjust over the following years, and what does that adjustment do to the revenue estimate? The Office for Budget Responsibility (OBR) answers the second question with a large-scale macroeconometric model in the Treasury/Cowles-Commission tradition---a model whose equation listing and variable definitions are published, but which is operated in proprietary EViews software, on a databank the OBR does not fully publish, and with judgement imposed through unpublished residual adjustments \citep{obr2013model, obr2025modelpage}.

This paper describes an open-source Python emulator of that model, built as the near-term fiscal-scoring member of the \emph{PolicyEngine Macro} suite.\footnote{\url{https://macromod.vercel.app/obr/}; source at \url{https://github.com/PolicyEngine/obr-macroeconomic-model}. PolicyEngine Macro is a suite of open macroeconomic simulation models---overlapping-generations, macroeconometric, and structural-VAR classes---that score the same PolicyEngine reform objects and report comparable real-world aggregates \citep{policyengine}.} The emulator takes the OBR's published model code of 15 October 2025---372 equations---transpiles it mechanically into Python, initialises it from the November 2025 \emph{Economic and Fiscal Outlook} (EFO) detailed forecast tables \citep{obr2025efo} and a 348-series snapshot of Office for National Statistics data, and solves it by Gauss--Seidel iteration. Policy analysis proceeds the way the OBR's own does: the baseline is anchored to the published forecast through add-factors, a shock is applied to an exogenous lever (or injected through a microsimulation bridge), and the reform effect is read as the deviation between the shocked and baseline solutions.

The contribution is threefold. First, \emph{replication as infrastructure}: to my knowledge this is the first publicly runnable implementation of the OBR's macroeconomic model. The OBR publishes the equation listing and a briefing paper describing the model \citep{obr2013model} but no official working paper accompanies the model itself, and the model cannot be run without EViews and the OBR's internal databank. The emulator makes the official forecaster's transmission mechanisms---the marginal propensities, error-correction speeds, and accounting identities that shape UK fiscal scoring---inspectable and executable by anyone with Python.

Second, \emph{integration with microsimulation}: a static-costing bridge connects PolicyEngine's UK tax--benefit microsimulation \citep{policyengine} to the macro model, reproducing the OBR's own policy-costing workflow \citep{obr2014costings} in which a static costing from the revenue departments enters the macroeconomic forecast and second-round demand effects come out. A reform expressed as PolicyEngine parameters is costed at household resolution, converted to a shock on nominal household disposable income, and propagated through the model's consumption and expenditure blocks.

Third, \emph{honest validation}: the project distinguishes sharply between what the model reproduces by construction and what it genuinely forecasts. The anchored baseline matches the EFO to within 0.2--0.3 per cent by design and this invariant is hard-gated in continuous integration; the free-running model, scored without the add-factors and de-seeded from the EFO path, is documented as weak---only 4 of 11 computed headline variables land within band---with each residual error traced either to a fixed bug or to a documented ``floor'' imposed by constants the OBR does not publish. Section~\ref{sec:limitations} places the model on the suite's \emph{verification gradient}, which grades each model class by how far its outputs can be audited against ground truth.

The rest of the paper proceeds as follows. Section~\ref{sec:literature} places the project in the UK macroeconometric modelling tradition and the replication literature. Section~\ref{sec:model} sets out the model's equation structure block by block---the ECM template, the expenditure core, consumption, investment, the labour market, trade, government spending, incomes, prices, and the full receipts-to-borrowing fiscal chain---together with the add-factor architecture and the anchoring mechanism. Section~\ref{sec:method} describes the implementation: transpiler, solver, closure swap, microsimulation bridge, and hosted interfaces. Section~\ref{sec:calibration} presents calibration and validation against the November 2025 EFO. Section~\ref{sec:results} reports a worked reform---1p on the basic rate of income tax---and compares it with HMRC's published ready reckoner \citep{hmrc2025reckoner}; Section~\ref{sec:experiments} adds three further experiments through the direct macro levers (a government-consumption expansion, a VAT-style price shock, and a corporation-tax rise). Section~\ref{sec:limitations} discusses limitations, and Section~\ref{sec:conclusion} concludes. Appendices provide a variable glossary, the complete free-running scorecard, and solver details.
